\begingroup
\renewcommand{\arraystretch}{1.45}
\clearpage\fancyhead[L]{\small TRANSZENDENTE FUNKTIONEN}
\pdfbookmark[0]{TEIL II. TRANSZENDENTE FUNKTIONEN}{agm-pars2}
\SourceLeaf{381}
\begin{center}
\textbf{TEIL II.}\par
\textbf{ÜBER DIE TRANSZENDENTEN FUNKTIONEN, DIE AUS DER DIFFERENTIATION}\par
\textbf{ARITHMETISCH-GEOMETRISCHER MITTEL ENTSTEHEN.}
\end{center}
\par\begin{center}\rule{.15\linewidth}{.4pt}\end{center}
\Needspace{43\baselineskip}
\AGMArticle{9}
Es seien
\[
\left\{\begin{array}{llllll}
x,&x',&x'',&x'''\; .\; .\; .\; .&&\text{(I)}\\[6pt]
y,&y',&y'',&y'''\; .\; .\; .\; .&&\text{(II)}
\end{array}\right.
\]
Folgen, die ebenso gebildet sind wie die Folgen in Artikel 1, nämlich so, dass jeder Term in
$\left\{\begin{array}{c}\mathrm{I}\\\mathrm{II}\end{array}\right\}$ das
$\left\{\begin{array}{c}\text{arithmetische}\\\text{geometrische}\end{array}\right\}$
Mittel zwischen den beiden vorhergehenden Termen ist, und somit $x^{\infty}=y^{\infty}=\mathrm{M}(x,y)$:
Offenbar sind alle diese Größen und daher auch ihr Grenzwert Funktionen
der beiden Variablen $x,y$ und ändern sich, sobald die eine oder beide
eine Änderung erfahren. Hier werden wir nur von unendlich kleinen Änderungen sprechen.
\Needspace{24\baselineskip}
Es ergibt sich daher
\[
\begin{aligned}&\mathrm{d}x'=\tfrac12\mathrm{d}x+\tfrac12\mathrm{d}y,\\&
\mathrm{d}y'=\tfrac12\surd\frac{y}{x}. \mathrm{d}x
+\tfrac12\surd\frac{x}{y}. \mathrm{d}y\\&=\tfrac12\frac{y'}{x}\mathrm{d}x+\tfrac12\frac{y'}{y}\mathrm{d}y\end{aligned}
\]
und. ebenso
\[
\begin{aligned}&\mathrm{d}x''=\tfrac12\mathrm{d}x'+\tfrac12\mathrm{d}y',\\&
\mathrm{d}y''=\tfrac12\frac{y''}{x'}. \mathrm{d}x'
+\tfrac12\frac{y''}{y'}\mathrm{d}y'\quad\text{usw.}\end{aligned}
\]
Auf diese Weise könnten die Differentiale aller Terme in I und II durch die Differentiale von
$x,y$ mithilfe von Substitutionen ausgedrückt werden; das Gesetz der Folge würde dabei aber in sehr
unübersichtlicher Form hervortreten. Um es uns möglichst deutlich vor Augen zu führen, werden wir die folgenden
Abkürzungen verwenden: Wir werden schreiben
\Needspace{20\baselineskip}
\SourceLeaf{382}
\[
\begin{array}{lll}
\dfrac{\mathrm{d}x}{x}+\dfrac{\mathrm{d}y}{y}=f,&
\dfrac{\mathrm{d}x'}{x'}+\dfrac{\mathrm{d}y'}{y'}=f',&
\dfrac{\mathrm{d}x''}{x''}+\dfrac{\mathrm{d}y''}{y''}=f''\quad\mathrm{usw.}\\[7pt]
\dfrac{\mathrm{d}x}{x}-\dfrac{\mathrm{d}y}{y}=g,&
\dfrac{\mathrm{d}x'}{x'}-\dfrac{\mathrm{d}y'}{y'}=g',&
\dfrac{\mathrm{d}x''}{x''}-\dfrac{\mathrm{d}y''}{y''}=g''\quad\mathrm{usw.}
\end{array}
\]
\Needspace{34\baselineskip}
woraus sich sofort ergibt
\[
\begin{aligned}
f'&=\tfrac12\left(\frac{\mathrm{d}x}{x'}+\frac{\mathrm{d}x}{x}+\frac{\mathrm{d}y}{x'}+\frac{\mathrm{d}y}{y}\right)\\&=f+\tfrac12\frac{\mathrm{d}x}{xx'}(x-x')+\tfrac12\frac{\mathrm{d}y}{yx'}(y-x')\\[7pt]
&=f+\tfrac12\frac{\mathrm{d}x}{xx'}(x-x')-\tfrac12\frac{\mathrm{d}y}{yx'}(x-x')\\&=f+\tfrac12 g\frac{x-x'}{x'}
\end{aligned}
\]
und auf ganz ähnliche Weise
\[
f''=f'+\tfrac12 g'\frac{x'-x''}{x''},\quad
f'''=f''+\tfrac12 g''\frac{x''-x'''}{x'''}\quad\mathrm{usw.}
\]
\Needspace{27\baselineskip}
Ebenso ergibt sich
\[
\begin{aligned}&g'=\tfrac12\left(\frac{\mathrm{d}x}{x'}-\frac{\mathrm{d}x}{x}+\frac{\mathrm{d}y}{x'}-\frac{\mathrm{d}y}{y}\right)\\&=\tfrac12\frac{\mathrm{d}x}{xx'}(x-x')+\tfrac12\frac{\mathrm{d}y}{yx'}(y-x')\\&=\tfrac12 g\left(\frac{x-x'}{x'}\right)\end{aligned}
\]
und auf dieselbe Weise \(g''=\tfrac12 g'\frac{x'-x''}{x''},\quad
g'''=\tfrac12 g''\frac{x''-x'''}{x'''}\) usw.
Daraus ist auch ersichtlich, dass
\(f'=f+g',\ f''=f+g'+g'',\ f'''=f+g'+g''+g'''\) usw.
\Needspace{25\baselineskip}
Ohne Mühe erkennt man, dass die Folge
\(\frac{g'}{g},\frac{g''}{g},\frac{g'''}{g}\) usw.
äußerst schnell unter jede Grenze abnimmt; daher werden wir durch eine äußerst schnell konvergierende unendliche Reihe erhalten
\[
\begin{aligned}&f^\infty=f+g\Bigl\{\tfrac12\frac{x-x'}{x'}
+\tfrac14\frac{x-x'}{x'}\cdot\frac{x'-x''}{x''}\\&\qquad+\tfrac18\frac{x-x'}{x'}\cdot\frac{x'-x''}{x''}\cdot\frac{x''-x'''}{x'''}
+\mathrm{usw.}\Bigr\},\\&g^\infty=0\end{aligned}
\]\Needspace{24\baselineskip}

daher
\[
\begin{aligned}
\frac{2\mathrm{d}x^\infty}{x^\infty}
&=\frac{2\mathrm{d}y^\infty}{y^\infty}
=\frac{2\mathrm{d}\mathrm{M}(x,y)}{\mathrm{M}(x,y)}\\[7pt]
&=\frac{\mathrm{d}x}{x}+\frac{\mathrm{d}y}{y}\\&\quad+\left\{\frac{\mathrm{d}x}{x}-\frac{\mathrm{d}y}{y}\right\}\times
\left\{\tfrac12\frac{x-x'}{x}+\tfrac14\frac{x-x'}{x'}\cdot\frac{x'-x''}{x''}+\mathrm{usw.}\right\}
\end{aligned}
\]
\Needspace{28\baselineskip}
oder
\[
\begin{aligned}
&\mathrm{d}\mathrm{M}(x,y)=\mathrm{M}(x,y)\times\Bigl\{\\&\quad\frac{\mathrm{d}x}{x}\times\left(\tfrac12+\tfrac14\frac{x-x'}{x'}+\tfrac18\frac{x-x'}{x'}\cdot\frac{x'-x''}{x''}+\mathrm{usw.}\right)\\&\quad+\frac{\mathrm{d}y}{y}\times\left(\tfrac12-\tfrac14\frac{x-x'}{x'}-\tfrac18\frac{x-x'}{x'}\cdot\frac{x'-x''}{x''}-\mathrm{usw.}\right)\Bigr\}
\end{aligned}
\]
\Needspace{39\baselineskip}
\par Eine etwas bequemere Rechnung erhalten wir auf folgende Weise: Es ergibt sich
\[
\begin{aligned}&\frac{x-x'}{x'}=\frac{x-y}{x+y}=\frac{(x-y)^2}{xx-yy}\\&=\frac{(x+y)^2-4xy}{xx-yy}=4\frac{x'x'-y'y'}{xx-yy}\text{\textsuperscript{*)}}\end{aligned}
\]\footnotetext{\csname HistoricalNoteText30\endcsname}
woraus folgt
\[
\begin{aligned}
&\mathrm{d}\mathrm{M}(x,y)=\frac{\mathrm{M}(x,y)}{2(xx-yy)}\times\Bigl\{\\&\quad\frac{\mathrm{d}x}{x}\times\bigl(xx-yy+2(x'x'-y'y')\\&\qquad+4(x''x''-y''y'')+8(x'''x'''-y'''y''').\;.\;.\bigr)\\&\quad+\frac{\mathrm{d}y}{y}\times\bigl(xx-yy-2(x'x'-y'y')\\&\qquad-4(x''x''-y''y'')-8(x'''x'''-y'''y''').\;.\;.\bigr)\Bigr\}
\end{aligned}
\]
\Needspace{26\baselineskip}
\SourceLeaf{383}
\par 
Damit sich die Rechnung mithilfe von Logarithmen möglichst leicht ausführen lässt, fügen wir die folgenden Formeln hinzu, die sich unmittelbar aus dem Algorithmus der Folgen I. II ergeben:
\[
(x'x'-y'y')=\frac{(xx-yy)^2}{16x'x'},\quad
x''x''-y''y''=\frac{(x'x'-y'y')^2}{16x''x''}\quad\mathrm{usw.}
\]
Diese zeigen zugleich sofort, wie schnell unsere Reihen konvergieren müssen.

\par Wenn es mehr zusagt, können auch die folgenden Formeln verwendet werden:
\[
x'x'-y'y'=\tfrac14(x-y)^2,\quad
x''x''-y''y''=\tfrac14(x'-y')^2\quad\mathrm{usw.}
\]

\clearpage
\endgroup
